%=============================================================================!
% 8.7  BONDED FORCE FIELDS                                                    !
%=============================================================================!
\section{Bonded force fields}
\label{sec:pe-bonded}

The solvents of a lithium-ion battery's electrolyte are small organic
molecules, carbonates such as ethylene carbonate, a ring of five atoms,
and chains such as dimethyl carbonate. While no reaction takes place, each
atom of a molecule keeps the same bonded partners, and its energy can be
written as a sum of terms for the bonds, the angles between bonds and the
twists about them, with the interactions of atoms farther apart, in the
same molecule or in different ones, added as pairs.
This section sets out that form, called a force field, and the one new
piece of geometry it needs, the dihedral angle.

\subsection*{The terms}

In the all-atom form, in which every hydrogen is a site of its own, of
the force field OPLS (Optimized Potentials for Liquid Simulations) of
Jorgensen and co-workers\cite{Jorgensen1996},
\begin{equation}
   \begin{aligned}
   U &= \sum_{\text{bonds}}\kappa_r(r - r_{\mathrm{eq}})^2
   + \sum_{\text{angles}}\kappa_\theta(\theta - \theta_{\mathrm{eq}})^2
   + \sum_{\text{dihedrals}}\sum_{n=1}^{3}\frac{\kappa_n}{2}\bigl[1 + (-1)^{n+1}
   \cos n\psi\bigr]\\
   &\quad + \sum_{i<j}f_{ij}\Bigl[\frac{k_{\mathrm e}q_iq_j}{r_{ij}}
   + 4\varepsilon_{ij}\Bigl(\frac{\sigma_{ij}^{12}}{r_{ij}^{12}}
   - \frac{\sigma_{ij}^6}{r_{ij}^6}\Bigr)\Bigr] .
   \end{aligned}
   \label{eq:pe-opls}
\end{equation}
The first sum holds each bond near its length $r_{\mathrm{eq}}$ by a
spring; $\kappa_r$ is half the stiffness $k$ of the spring energy $\frac12
kx^2$ (\cref{sec:en-varying}), since the factor $\frac12$ is absorbed into
it. The second holds each angle
$\theta$ between two bonds at an atom near $\theta_{\mathrm{eq}}$, as the
three atoms of \cref{fig:ha-noether}(a) were held, with $\kappa_\theta$
half the angular stiffness for the same reason. The third gives the
energy of a twist $\psi$ about each bond as a short series of cosines,
written out $\frac12\kappa_1(1 + \cos\psi) + \frac12\kappa_2(1 - \cos2\psi)
+ \frac12\kappa_3(1 + \cos3\psi)$ with constants $\kappa_n$, and
the last adds the Coulomb energy of charges $q_i$ placed on the atoms
(\cref{sec:pe-coulomb}) and a Lennard-Jones energy, for every pair. The
factor $f_{ij}$ is 0 for atoms joined by one bond or two, whose
interaction the bond and angle terms already describe, $\frac12$ for
atoms three bonds apart, and 1 for all others, including every pair in
different molecules. Each constant depends on the kinds of atom involved,
and a force field is a table of them, fitted to calculations of the
electrons and to measured properties of liquids.

\subsection*{The dihedral angle}

Four atoms bonded in a chain, $a$--$b$--$c$--$d$, can twist about the
middle bond $b$--$c$ (\cref{fig:pe-dihedral}). The twist is measured by
the angle between the plane of $a$, $b$, $c$ and the plane of $b$, $c$,
$d$, seen along the middle bond. Each plane is fixed by the arrow at right
angles to it, which the cross product of two arrows in the plane gives
(\cref{sec:ro-cross}). With $\vb{b}_1 = \vb{r}_b - \vb{r}_a$, $\vb{b}_2 =
\vb{r}_c - \vb{r}_b$ and $\vb{b}_3 = \vb{r}_d - \vb{r}_c$ along the
three bonds,
\begin{equation*}
   \vb{n}_1 = \vb{b}_1\times\vb{b}_2 , \qquad
   \vb{n}_2 = \vb{b}_2\times\vb{b}_3
\end{equation*}
are at right angles to the first and second planes. Seen along
$\vb{b}_2$, the two planes appear as lines from the middle bond towards
$a$ and towards $d$, and $\psi$ is the angle between them. The part of
$\vb{b}_1$ along $\vb{b}_2$ drops out of $\vb{n}_1$, since $\vb{b}_2\times
\vb{b}_2 = \vb{0}$, and the part left points from $a$ towards the bond, so
$\vb{n}_1$ is $\vb{b}_2$ crossed with an arrow from the bond towards $a$;
likewise $\vb{n}_2$ is $\vb{b}_2$ crossed with an arrow from the bond
towards $d$. Each is its line turned through a right angle about
$\vb{b}_2$, the same way for both, so the angle between them is $\psi$,
and by the dot product
\begin{equation}
   \cos\psi = \frac{\vb{n}_1\cdot\vb{n}_2}{\lvert\vb{n}_1\rvert\,
   \lvert\vb{n}_2\rvert} .
   \label{eq:pe-dihedral}
\end{equation}
The cosine does not tell a twist one way from the same twist the other
way. The cross product $\vb{n}_1\times\vb{n}_2$ is along $\vb{b}_2$, the
one direction at right angles to both, pointing along $\vb{b}_2$ or
against it according to the sense of the twist, so the sign of
$(\vb{n}_1\times\vb{n}_2)\cdot\vb{b}_2$ gives the sign of $\psi$, which
then runs from $-\pi$ to $\pi$; $\psi$ is positive when, seen from $b$
looking towards $c$, the bond to $a$ turns clockwise to reach the bond to
$d$. For a flat zigzag, $a$ and $d$ on
opposite sides of the middle bond, $\vb{n}_1$ and $\vb{n}_2$ point in
opposite directions and $\psi = \ang{180}$, the arrangement called
trans; with $a$ and $d$ on the same side, $\psi = 0$, called cis. The
function \texttt{dihedral\_angle} of \texttt{mdlab} computes $\psi$ in
this way.

\begin{figure}[htb]
\centering
\begin{tikzpicture}[line cap=round, line join=round, scale=1.2,
   atom/.style={circle, fill=black, inner sep=1.6pt},
   arr/.style={-{Stealth[length=5pt]}, line width=0.8pt, accent}]
   \draw[line width=1pt] (0,0) -- (90:1.3) node[atom, label={above:$a$}] {};
   \draw[line width=1pt] (0,0) -- (-30:1.3) node[atom, label={right:$d$}] {};
   \node[atom, label={left:{$b$ in front of $c$}}] at (0,0) {};
   \draw[arr] (0,0) -- (0:1.0) node[right] {$\vb{n}_1$};
   \draw[arr] (0,0) -- (240:1.0) node[below] {$\vb{n}_2$};
   \draw[black!50] (90:0.5) arc[start angle=90, end angle=-30, radius=0.5];
   \node at (45:0.72) {$\psi$};
\end{tikzpicture}
\caption{Four atoms bonded in a chain, seen along the middle bond with
$b$ in front of $c$. The planes of $a$, $b$, $c$ and of $b$, $c$, $d$ are
seen edge-on as the lines towards $a$ and $d$, at the dihedral angle
$\psi$; $\vb{n}_1$ and $\vb{n}_2$ are those lines turned through a right
angle the same way, so they meet at the same angle.}
\label{fig:pe-dihedral}
\end{figure}

\begin{keyidea}
A force field writes the energy of molecules as springs for bonds and
angles, a cosine series for each twist, and Coulomb and Lennard-Jones
energies between atoms more than two bonds apart. The twist is the
dihedral angle between the two planes of a chain of four atoms, found
from two cross products.
\end{keyidea}

\notebook{08}{twist a chain of four atoms and watch its torsion energy}

\begin{selfcheck}
In \cref{eq:pe-opls}, which values of $\psi$ make the $\kappa_1$ and the
$\kappa_3$ terms smallest, when both constants are positive?
\end{selfcheck}
